% Candidate D.74. Notes only, pending full independent audits.
\subsection{Integral Alexander cocycles at every finite ramification set}

\begin{statement}{D.74}{Proved}{Every finitely ramified Alexander first cohomology is an infinite countable product}
Let \(s=(a,b)\) be any finitely ramified abelian pair for the
tripod over \(\mathbb Q\), and let \(A_s=\Lambda C\), with
\(\Lambda=\mathbb Z_5[[u,v]]\), carry the actual action of D.35.
Then, as compact \(\mathbb Z_5\)-modules,
\[
 H^1_{\rm cts}(G_{\mathbb Q},A_s)
                  \simeq\prod_{\mathbb N}\mathbb Z_5.
                                                        \tag{D.74.1}
\]
The isomorphism is noncanonical. Thus every nonempty
geometrically pro-five metabelian section fibre above a
finitely ramified pair is a torsor under this compact group.
Every finitely ramified cusp pair has uncountably many
metabelian classes which are neither cuspidal nor induced
by rational points.

The conclusion concerns first cohomology and already nonempty
fibres. It does not assert existence above an arbitrary pair,
evaluate the remaining compact second-cohomology terms, or
compute a first higher-derived obstruction. In particular the
zero-rank twists of D.60--D.66 must be infinitely ramified.
\end{statement}

\begin{proof}
The proof constructs integral continuous cocycles before any
rank comparison. We distinguish the operator ring from its
actual module throughout.

\paragraph{A finite arithmetic presentation of positive deficiency.}
Choose a finite rational set \(S\) containing \(5,\infty\)
and the support of \(a,b\). Put
\[
 K=\mathbb Q(\zeta_5),\quad G=G_{K,S},\quad P=G^{(5)}.
\]
Let \(d=\dim_{\mathbb F_5}H^1(G,\mathbb F_5)\) and
\(h=\dim_{\mathbb F_5}H^2(G,\mathbb F_5)\).
These are finite, and the global Euler characteristic is
\[
                         1-d+h=-2,\qquad d-h=3.
                                                        \tag{D.74.2}
\]
Here \(K\) is totally imaginary of degree four and \(S\)
contains every place over five. This is the finite-module
global Euler characteristic in Milne ADT, Chapter I,
Theorem 5.1 \cite{DE_MilneADT}, with trivial
\(\mathbb F_5\) coefficients.
There is no real-place correction.

Write \(N=\ker(G\to P)\). As in D.64, an invariant
nonzero character \(N\to\mathbb F_5\) would give a
pro-five quotient of \(G\) with nontrivial image of \(N\).
Consequently \(H^1(N,\mathbb F_5)^P=0\), and the
Hochschild--Serre sequence gives
\[
 H^1(P,\mathbb F_5)=H^1(G,\mathbb F_5),\qquad
 H^2(P,\mathbb F_5)\lhook\joinrel\longrightarrow
                                     H^2(G,\mathbb F_5).
                                                        \tag{D.74.3}
\]
Thus \(P\) has a finite pro-five presentation with \(d\)
generators and \(r\le h\) relations. One can obtain it
directly from a minimal free presentation: its relation
coinvariants modulo fifth powers are dual to
\(H^2(P,\mathbb F_5)\); lifts of a finite basis normally
generate the kernel by the finite-five-group Nakayama
argument reproduced in D.64. In particular
\[
                              d-r\ge3.                 \tag{D.74.4}
\]
No assertion that \(P\) has cohomological dimension two
is used here.

The actual Alexander action factors through
\(\mathcal G=P\rtimes C_4\), a quotient of
\(G_{\mathbb Q,S}\). Indeed its finite coefficient
actions over \(K\) are five-groups, by the augmentation
and five-adic filtrations. The same holds for the finite
affine graphs of actual cocycles. Hence all first
cohomology in question factors through \(\mathcal G\).
The extension by \(C_4=\operatorname{Gal}(K/\mathbb Q)\)
has a complement by the coprime complement theorem,
or by compatible complements in its finite quotients.
Fix a generator \(c\) of such a complement. Its actual
cyclotomic multiplier is a primitive fourth root
\(\omega\in\mathbb Z_5^\times\).

\paragraph{The actual completed operator ring.}
Put \(\epsilon=1+5\), and let
\(\gamma\) act on \(\Lambda\) by the substitution
\[
 \varphi(t_i)=t_i^\epsilon,\qquad t_1=1+u,\ t_2=1+v.
\]
Use the completed group ring
\[
 R=\mathbb Z_5[[\,\mathbb Z_5^2\rtimes_\epsilon\mathbb Z_5\,]],
 \qquad z=\gamma-1.
                                                        \tag{D.74.5}
\]
It has the unique ordered expansion
\[
 R=\left\{\sum_{k\ge0}a_k(u,v)z^k:a_k\in\Lambda\right\},
 \qquad
 z a=\varphi(a)z+\varphi(a)-a.
                                                        \tag{D.74.6}
\]
This is the actual semidirect-product Iwasawa ring, as in
Venjakob, \emph{A noncommutative Weierstrass preparation
theorem}, Example 2.3 and its preceding construction
\cite{DE_Venjakob2002Skew}.
The additive expansion also follows from the ordered
coordinates of the three copies of \(\mathbb Z_5\).
The topology is the \((5,u,v,z)\)-adic topology.

Every \(p\in P\) acts on \(A_s\) as
\[
 \rho_s(p)=b_p\gamma^{e(p)},\qquad
 b_p\in\Lambda^\times,\quad b_p(0,0)=\chi_5(p)^2,\quad
 \chi_5(p)=\epsilon^{e(p)}.
                                                        \tag{D.74.7}
\]
Formula D.35.1 gives \(b_p\), including its actual Ihara
factor. Thus \(p\mapsto b_p\gamma^{e(p)}\) is a continuous
homomorphism \(P\to R^\times\), with image in
\(1+(5,u,v,z)R\). This is an equality of actual operators,
not an approximation by a weight.

The ring \(R\) acts on \(\Lambda\) by multiplication and
the indicated substitution. This completed action is
well-defined: \(\varphi-1\) increases the
\((5,u,v)\)-adic order by at least one, so the series
in (D.74.6) converges uniformly on \(\Lambda\).
The action is continuous. We prove below both its
faithfulness and the stronger infinite-rank property
needed in the argument.

We will also use that \(R\) is a left and right
Noetherian domain with a division ring of fractions \(D\).
Here is a filtration proof in this particular ring.
Give \(5\) weight two and \(u,v,z\) weight one.
The resulting complete filtration is equivalent to the
maximal-ideal topology. The relation
\[
 \varphi(u)-u=5u+15u^2+20u^3+15u^4+6u^5+u^6
\]
and its \(v\)-analogue increase weight by at least two.
Consequently all commutators of the three weight-one
symbols have weight strictly greater than two, and the
ordered expansion gives
\[
             \operatorname{gr}R
                   =\mathbb F_5[P_0,U_0,V_0,Z_0],
 \quad \deg P_0=2,\quad
 \deg U_0=\deg V_0=\deg Z_0=1.                         \tag{D.74.8}
\]
This is the weighted filtration, not the generally
noncommutative maximal-ideal associated graded ring.
The graded domain proves that \(R\) is a domain.
For a right ideal, lift a finite homogeneous generating
set of its associated graded ideal. Successive
subtraction raises the weight; completeness sums the
coefficients into finitely many elements of \(R\).
This proves finite generation of the original right
ideal. The identical argument works on the left.

For clarity, the Ore property follows without an
unproved flatness claim about \(\Lambda\). In a right
Noetherian domain two nonzero principal right ideals
cannot have zero intersection: if \(aR\cap bR=0\),
then \(bR,abR,a^2bR,\ldots\) form a direct sum in
\(R\), giving an infinite strictly ascending chain
of right ideals. Thus common nonzero right multiples
exist; likewise on the left. These Ore fractions give
the division ring \(D\), and localization is exact on
right modules by the common-denominator construction.
The general Noetherian assertion also follows from
Venjakob's Iwasawa structure theorem, Theorem 3.26
\cite{DE_Venjakob2001Structure},
whose hypotheses hold for this torsion-free uniform
pro-five group.

\paragraph{Fox cocycles with operator coefficients.}
Give \(R\) the left \(P\)-action in (D.74.7).
Evaluation on the chosen generators identifies its
continuous cocycles and coboundaries with
\[
 Z_R=\ker\{L:R^d\to R^r\},\qquad
 B_R=\{\,\mathbf r\,t:t\in R\,\},\quad
 \mathbf r=(\rho_s(p_i)-1)_i.                           \tag{D.74.9}
\]
Here \(L\) is the actual completed Fox relation matrix.
These are right \(R\)-modules. To verify continuity,
arbitrary values on free pro-five generators give
compatible finite affine five-group homomorphisms
at every finite ring quotient. The relator conditions
are exactly \(L\mathbf v=0\), and then the cocycle
descends through the closed normal relation subgroup.
Thus no abstract discontinuous cocycles enter (D.74.9).

By D.35 there is \(\delta\in P\), with
\(\chi_5(\delta)=1\), such that
\[
                 h=\rho_s(\delta)-1\in\Lambda
                                      \quad\hbox{is nonzero}.
                                                        \tag{D.74.10}
\]
Its existence uses the primitive actual Soul\'e element
of D.35, which cannot be canceled by the linear Kummer
multiplier. Its image in the finite-ramification
quotient still has the same nonzero operator.
Add \(\delta\) as a redundant presentation generator
and its defining relation. This increases \(d,r\)
equally; write \(i\) for this generator's index.
The \(i\)-th entry of \(\mathbf r\) is the
nonzero multiplication operator \(h\).

After right localization, the vector space
\[
             M_D=(Z_R/B_R)\otimes_R D
\]
has dimension at least \(d-r-1\ge2\). Indeed the
relation matrix has at most \(r\) independent rows,
and the nonzero coboundary column has rank one.
This is a statement about the operator coefficient
ring. We do not tensor the actual module
\(\Lambda\) with \(D\).

\paragraph{Tame semilinear descent and an invariant denominator.}
Conjugation by the actual operator \(\rho_s(c)\)
defines an order-four automorphism \(\alpha\) of \(R\).
For \(a\in\Lambda\) it is \(a(t_i^\omega)\);
on \(\gamma\) it is conjugation by the unit
multiplier \(b_c\), since the two substitutions commute.
It preserves the preceding topology and weighted
filtration. It acts on the leading symbols \(u,v\)
by \(\omega\) modulo five. It extends to \(D\).
Indeed \(\varphi(b_c)-b_c\) has weight at least
three, so \(\alpha(z)-z\) does also; the leading
symbols of \(5,z\) are unchanged.
Every nonidentity power of \(\alpha\) is outer on
\(D\): the valuation from (D.74.8) extends to Ore
fractions, and its associated graded division ring
is commutative. Therefore an inner automorphism acts
identically on leading symbols, whereas
\(\alpha^j\), \(1\le j\le3\), multiplies the leading
symbol of \(u\) by \(\overline\omega^j\ne1\).

The action on actual \(R\)-valued cocycles is
\[
          (c\cdot z)(p)=\alpha\bigl(z(c^{-1}pc)\bigr).
                                                        \tag{D.74.11}
\]
It preserves \(Z_R,B_R\), and gives \(M_D\) a
semilinear order-four action.

We use the following elementary independence fact.
Distinct automorphisms of a division ring, no two
differing by an inner automorphism, are linearly
independent as functions with coefficients on the
left. The same holds when they are tested only on
a subring that generates the division ring.
For a proof take a shortest nonzero relation.
Evaluate it at \(yx\), subtract a suitable left
multiple of its value at \(x\), and cancel one
term. Minimality then forces any two remaining
automorphisms to differ by conjugation by a
coefficient, a contradiction. Testing \(x,y\) in
the generating subring suffices, since the resulting
identity extends to all fractions.

Choose a nonzero class \(m\in M_D\) represented by
a cycle in \(Z_R\). There is \(t\in R\) for which
\(\sum_{j=0}^3c^j(m t)\ne0\). Otherwise, in one
nonzero coordinate of the right \(D\)-vector space
\(M_D\), this would be a nontrivial left-coefficient
relation among \(\alpha^j(t)\), contrary to the
preceding independence. Averaging the representing
cycle gives an exactly \(C_4\)-invariant
\[
                 \mathbf w\in Z_R,\qquad
                    \mathbf w\notin\mathbf r D.       \tag{D.74.12}
\]
Division by four is allowed in \(\mathbb Z_5\).

We may clear a denominator without losing this exact
invariance. Let
\[
 J=\{t\in R:h^{-1}w_i t\in R\}.
 \]
It is a nonzero right ideal by the Ore condition.
The intersection \(J_*=\bigcap_{j=0}^3\alpha^j(J)\)
is again nonzero, by successive common right multiples,
and is stable under \(\alpha\). Choose \(0\ne j_*\in J_*\).
The same independence argument, now in \(D\), shows
that for some \(t\in R\) the element
\[
                    d_*=\sum_{j=0}^3\alpha^j(j_*t)
 \]
is nonzero. It belongs to \(J_*\) and is
\(\alpha\)-invariant. Replace \(\mathbf w\) by
\(\mathbf w d_*\). It is still an exactly invariant
cycle with nonzero class in \(M_D\), and now
\[
             w_i=h u_*\quad\hbox{for some }u_*\in R.
                                                        \tag{D.74.13}
\]
Since it is not a division-ring coboundary, for some
generator index \(j\) the actual operator
\[
                  T=r_j u_*-w_j\in R
                               \quad\hbox{is nonzero}. \tag{D.74.14}
\]
This supplies the denominator control needed below;
an arbitrary Reynolds average was not assumed nonzero.

\paragraph{Every nonzero operator has infinite rank on the actual tame fixed module.}
Write \(\rho_s(c)=b_c\varphi_\omega\) on \(\Lambda\).
Its constant term is \(\omega^2\), and it has
order four. The integral eigenprojector
\[
       e=\frac14\sum_{j=0}^3\omega^{-2j}\rho_s(c)^j(1)
                                                        \tag{D.74.15}
\]
has constant term one, hence is a unit, and
\(\rho_s(c)e=\omega^2e\).
Also put
\[
 U=\frac14\sum_{j=0}^3\omega^{-j}\varphi_\omega^j(u),
 \qquad
 V=\frac14\sum_{j=0}^3\omega^{-j}\varphi_\omega^j(v).
                                                        \tag{D.74.16}
\]
They are integral formal coordinates with leading
terms \(u,v\), and satisfy
\(\varphi_\omega(U)=\omega U\),
\(\varphi_\omega(V)=\omega V\).
Thus \(eU^n\in\Lambda^{\rho_s(c)}\) whenever
\(n\equiv2\pmod4\).

Let \(0\ne T\in R\). The operator \(Te\) is nonzero,
since \(R\) is a domain. Expand it uniquely in
\(U,V,z\), and let \(a\) be its least total
\((U,V)\)-degree. Its degree-\(a\) coefficient has
the form
\[
                  \sum_{i+j=a}U^iV^j f_{ij}(z),
          \qquad f_{ij}\in\mathbb Z_5[[z]],
 \]
with at least one nonzero \(f_{ij}\). The substitution
\(\varphi\) has linear part
\(U\mapsto\epsilon U,\ V\mapsto\epsilon V\).
Consequently the degree-\(n+a\) term of
\((Te)(U^n)\) is exactly
\[
       \sum_{i+j=a}f_{ij}(\epsilon^n-1)U^{n+i}V^j.
                                                        \tag{D.74.17}
\]
This identity includes the infinitely many completed
\(z\)-coefficients: they converge because
\(\epsilon^n-1\in5\mathbb Z_5\).

A nonzero series in \(\mathbb Z_5[[z]]\) has only
finitely many zeros in \(5\mathbb Z_5\), by applying
the elementary Strassmann theorem to \(f(5z)\).
Here this elementary form needs no hypothesis about
a rational formal series: \(f(5z)\) has coefficients
tending to zero. If \(N\) is the largest index at
which their maximal absolute value is attained,
division by \(z-a\), for a zero \(a\in\mathbb Z_5\),
is obtained by the convergent coefficient sums
\(h_i=\sum_{j>i}g_j a^{j-i-1}\). The resulting
strictly convergent quotient has largest maximal
index \(N-1\). Repeating proves that there are
at most \(N\) zeros, including multiplicity.
The values \(\epsilon^n-1\) are distinct. Thus for
all but finitely many positive \(n\equiv2\pmod4\),
(D.74.17) is nonzero. The corresponding actual
outputs \(T(eU^n)\) have different
\((U,V)\)-orders \(n+a\), and are linearly
independent over \(\mathbb Q_5\). In particular
\[
 \operatorname{rank}_{\mathbb Z_5}
          T\bigl(\Lambda^{\rho_s(c)}\bigr)=\infty.
                                                        \tag{D.74.18}
\]
The same argument without restricting \(n\) proves
faithfulness of the whole actual completed action.
There is no use of a purely formal operator field
acting faithfully by assumption.

\paragraph{Integral cocycles and their nonzero descended classes.}
For \(f\in\Lambda^{\rho_s(c)}\), evaluate every value
of the invariant \(R\)-valued cocycle \(\mathbf w\)
on \(f\). This gives an actual integral continuous
\(A_s\)-valued cocycle \(z_f\) on \(P\).
Equation (D.74.11) and \(c f=f\) make it exactly
\(C_4\)-equivariant. Extending it by zero on the
chosen complement gives a cocycle on
\(\mathcal G=P\rtimes C_4\), hence on
\(G_{\mathbb Q,S}\).

If its class is a coboundary, its values on \(P\)
have the form \(r_j m\) for one \(m\in\Lambda\).
The special component (D.74.13) gives
\[
                   h m=w_i f=h(u_*f).
 \]
Multiplication by the nonzero element \(h\) is
injective on the integral domain \(\Lambda\), so
\(m=u_*f\). Its \(j\)-th component now implies
\(T(f)=0\), where \(T\) is (D.74.14).
Therefore the kernel of the class map
\[
 \Lambda^{\rho_s(c)}\longrightarrow
                  H^1(G_{\mathbb Q,S},A_s),\quad
                       f\longmapsto[z_f]
 \]
is contained in the kernel of \(T\). By (D.74.18)
the image, and hence the cohomology group, has
infinite integral rank. Every step took place in
integral actual cochains.

Finally D.47 identifies this first cohomology with
the unrestricted group: outside \(S\) every integral
cocycle has zero inertia, by the tame positive-weight
argument already proved there. D.35 proves
torsion-freeness, and D.47 proves compactness and
that it is a product of at most countably many
copies of \(\mathbb Z_5\). Infinite rank therefore
gives precisely (D.74.1).

For a fixed pair each of the three cusp packets
contributes at most one metabelian class, by its
primitive abelian boundary coordinate and the
packet formulas of D.35. Rational points form a
countable set. A nonempty torsor under
\(\prod_{\mathbb N}\mathbb Z_5\) is uncountable,
so removing these classes still leaves uncountably
many. A finitely ramified cusp pair has an actual
cusp section, so its fibre is nonempty. No assertion
of existence is made for other pairs.
\end{proof}
